\section{Build a digest from a sequence of bytes}
A SHA-256 digest is a sequence of 32 bytes computed from a message. In this project,
the message is initially an array in simulator memory. Later it will be the contents
of a file. A comparison program can calculate a digest for the same bytes and check
our result. The filename, the way the text is displayed, and the size of a read buffer
are outside the message.

The first test message is the three ASCII bytes \texttt{61 62 63}, which spell
\texttt{abc}. There is no trailing zero or newline in that message. We will follow
those bytes into a 64-byte block, decode the block into words, and observe the first
compression round. Once that one-block path is working, we will handle messages that
fill several blocks and messages that arrive in pieces.

A complete hash involves several responsibilities. The compression function changes
an eight-word state using one block. The streaming interface retains bytes that do
not yet make a full block. Finalization adds the message terminator and length, then
serializes the result. The file wrapper supplies bytes to that interface. Keeping these
boundaries visible gives us more useful observations than a single final mismatch.

For example, suppose \texttt{abc} produces an incorrect digest. If its first decoded
word is not \texttt{61626380}, inspect padding and byte order before studying the
round loop. If the schedule and first round agree but a 56-byte message fails, inspect
the finalization boundary. Each observation narrows the part of the program we need
to understand.
